\section{Virtual Machine Runtime}
\label{sec:runtime}

\subsection{Generated freestanding runtime}

The protector generates a complete Rust source file for each protection run. It
contains the shared ISA, cipher, and derivation modules; the engine; the encrypted
payload and decode table; the seed, expected hash, key shares, entry states, and
fuel; one \code{extern "C"} entry thunk per protected function; and safe
\texttt{\#[used]} function-pointer statics that keep the thunks alive through
linking without \code{no_mangle}.

The source is compiled with:

\begin{center}
\small
\begin{tabular}{@{}ll@{}}
\toprule
Setting & Value \\
\midrule
Edition & 2021 \\
Crate type & binary, \code{no_std}, \code{no_main} \\
Unsafe policy & \texttt{\#![forbid(unsafe\_code)]} in the default tier \\
Panic strategy & abort \\
Codegen units & 1 \\
Relocation model & PIC \\
Overflow checks & off (wrapping semantics are explicit) \\
Optimization & \code{opt-level=z} \\
Link mode & \verb|-nostdlib -static -no-pie|, custom linker script \\
Garbage collection & disabled (\verb|--no-gc-sections|) \\
Output segments & one RX \code{PT_LOAD} \\
\bottomrule
\end{tabular}
\end{center}

\code{opt-level=z} also avoids libc \code{memcpy} synthesis for the current
runtime. The linker script discards \code{.data}, \code{.bss}, \code{.got},
relocations, exception tables, comments, and notes. The blob validator rejects
undefined symbols, relocations, writable load segments, missing symbols, wrong
payload/seed/hash values, and entries outside the segment before injection.

\subsection{VM input and state}

The generated entry receives the twelve integer arguments permitted by the ABI
shim and passes the first six register arguments plus up to six stack arguments
to the engine. The engine input contains the encrypted payload, private data
blob, seed, expected hash, key shares, entry-state table, and fuel budget.

The runtime state is:

\begin{itemize}
  \item 96 virtual 64-bit registers, masked at rest;
  \item a 4 KiB private byte-addressable stack, masked at rest;
  \item a read-only data blob of at most 2 KiB;
  \item a current state identifier and a guest call stack;
  \item a fuel counter and a step counter;
  \item the decode tables, derived once per invocation;
  \item the effective payload key, derived once per invocation.
\end{itemize}

There is no mutable global engine state. The payload, tables, and seed are
read-only after initialization. This is the basis for thread and signal
reentrancy.

\subsection{Dispatcher}

The normal dispatcher performs the following loop:

\begin{enumerate}
  \item Derive the decode tables from the seed.
  \item Derive the effective key from the payload hash and key shares.
  \item Initialize virtual registers, private stack, data blob, and fuel.
  \item Set the current state to the requested entry state.
  \item Scan the encrypted manifest for the current state.
  \item Decode instructions sequentially from the block body.
  \item Execute each semantic operation, updating registers, flags, memory, or
        control state.
  \item On a terminator, update the state or return/trap; otherwise continue.
  \item Halt with \code{BadPayload} on any validation failure, or
        \code{FuelExhausted} when the budget is exhausted.
\end{enumerate}

Every manifest scan and every decoded instruction charges fuel. The default
budget is 2,000,000 steps. The budget is a deterministic bound, not a wall-clock
timeout; it prevents a malformed or adversarial payload from causing an
unbounded loop.

\subsection{Semantic operations}

The ISA has 48 semantics. They cover moves, immediate loads, wrapping
addition/subtraction/multiplication, bitwise operations, shifts, rotates,
comparisons, select, memory loads and stores at 8/16/32/64 bits, control-flow
transitions, return, trap, direct call, repeated moves/fills, and byte swap.
Appendix~\ref{app:isa} gives the canonical operands and semantics.

All arithmetic is 64-bit wrapping. Shift and rotate counts are masked to 63.
There is no division, no floating point, and no undefined shift. A memory
operation computes an effective address and checks the full width before
accessing the private stack. Invalid registers, invalid opcodes, malformed
operands, and out-of-range addresses return \code{None}, which the dispatcher
converts to \code{BadPayload}.

\subsection{Direct calls and the guest stack}

A selected function may call another selected function with a direct
\code{call rel32}. The lifter resolves the target to an entry state and emits a
\code{CALL} semantic. At runtime the VM pushes the return state on a bounded guest
call stack, switches to the callee entry, and returns to the saved state on
\code{RET}. The call stack is part of the private state; no host stack frame is
created for the guest call. A call to an external or unresolved target is a build
error, and an invalid return state is a \code{BadPayload} halt.

The compiler can automatically lift resolvable direct callees with
\verb|--follow-calls|. This is convenient for small helper clusters, but it does
not extend the ABI boundary: a callee that uses an unsupported instruction or an
external call still fails the build.

\subsection{Failure semantics}

The five runtime outcomes are:

\begin{description}[leftmargin=1.2em,style=nextline]
  \item[\code{Returned}] The program executed \code{RET}; the value is the
        virtual \code{RAX} contents.
  \item[\code{BadPayload}] A manifest, body, opcode, register, address, flag, or
        call-state validation failed. Value is zero.
  \item[\code{FuelExhausted}] The deterministic step budget was exhausted.
        Value is zero.
  \item[\code{Trap}] The program executed a \code{TRAP} semantic. Value is
        zero.
  \item[\code{LoopBound}] The fixed-pass constant-time frontier did not
        converge within the configured pass budget. Value is zero.
\end{description}

No failure path unwinds, allocates, or touches host memory. A failure is
deterministic for a fixed payload and input, which makes it testable and prevents
a malformed payload from turning into an environment-dependent crash.

\subsection{Memory safety argument}

The runtime is safe Rust. The public path uses checked slice access; there is no
\code{unsafe}, no raw pointer, no unchecked array indexing, no allocation, and no
panic in normal execution. The generated panic handler is a final \code{no_std}
guard and is not reachable through the supported semantic operations.

This is a source-level safety argument, not a proof of the entire interpreter.
Kani proves bounded memory helpers, Miri checks the Rust abstract machine on the
library tests, and libFuzzer explores malformed payloads. The evaluation section
states the exact scope of each result.
